\section{Conclusion}

% \pscomment{To Do: rewrite a much shorter version.}
% This paper develops a task-based framework for understanding how advances in AI reshape automation, the division of labor, firms' organizational structure, and production more generally.
% By modeling production as a sequence of steps that can be aggregated into tasks and then into jobs, we show how firms optimally allocate work between humans and AI (whether in augmented or fully automated form) when doing so requires trading off gains from specialization against coordination frictions.
% A central feature of the model is the possibility of AI chains, in which multiple steps are executed by AI and only the final output is verified by a human.
% Chaining can overturn step-level comparative advantage logic in factor assignment and non-linearly amplify the effects of marginal improvements in AI quality.

% We characterize short-run AI deployment decisions when job boundaries, and therefore the set of tasks assigned to each worker, are fixed, and the long-run joint optimization of job design and AI deployment strategy.
% We show that improvements in AI quality can induce discrete reorganizations of work.
% A key implication is that marginal increases in AI quality may generate little or no cost savings until a threshold is crossed, after which the optimal production structure changes discontinuously.
% This helps explain why firms invest so heavily in AI capabilities even when short run returns appear limited, and is consistent with the J-curve pattern of technology adoption \citep{brynjolfsson2021productivity} in which early adoption raises costs before the anticipated gains from reorganization are realized.

% We also show how the production functions derived from firm-level cost minimization can be mapped into a CES production function at the aggregate level when firms differ in how they deploy a common, general-purpose AI technology.
% Our framework therefore also allows studying the aggregate productivity and labor demand implications of improvements in AI quality through a micro-foundation for how individual firms respond to such advances.

% We complement the theoretical analysis with empirical evidence consistent with three core mechanisms implied by the model.
% Using a task-level dataset of AI exposure and AI execution outcomes, we find that AI-executed steps tend to appear in contiguous AI chains, that occupations with more fragmented AI-exposed steps in their workflow exhibit lower realized AI execution conditional on exposure levels, and that when comparable steps appear across occupations those adjacent to AI-executed neighbors are significantly more likely to be executed by AI themselves.
% These patterns align with the central economic forces highlighted by the model: the importance of positional relationships among steps, the role of AI-able step dispersion in shaping execution outcomes, and the local complementarities created by AI chains.




This paper develops a task-based framework for how advances in AI reshape automation, the division of labor, and firms' organizational structure. 
We model production as a fixed ordered sequence of steps that firms aggregate into tasks and then into jobs. 
Firms assign steps to humans and AI (in either augmented or automated form) while trading off specialization gains against coordination costs. 
A central feature of the model is AI chains, in which AI executes multiple steps jointly and a human verifies only the final output. 
Chaining can overturn step-level comparative advantage and generate non-linear gains from marginal improvements in AI quality. 

We characterize short-run AI deployment when job boundaries are fixed and long-run optimal execution modes of steps when firms jointly choose job design and AI deployment strategy. 
In the short run, the clustering or dispersion of AI-able steps, captured by our fragmentation index, determines the extent of gains from AI automation at the occupation level. 
In the long run, the framework clarifies how AI reshapes worker skill requirements through job redesign and provides a micro-foundation for J-curve adoption dynamics, with modest gains at low AI quality levels and sharply higher returns once workflow reorganization becomes optimal. 

We complement the theory with empirical evidence on three implications of the model. 
Using data on AI exposure and execution of O*NET tasks, we show that (1) AI-executed steps appear together in contiguous chains, (2) occupations with more fragmented AI-exposed steps exhibit lower realized execution conditional on exposure, and (3) when comparable steps appear across occupations, steps adjacent to AI-executed neighbors are more likely to be executed by AI themselves. 
Together, these findings highlight how positioning of steps in the workflow shapes local complementarities for AI and, in turn, its economic effects and adoption patterns. 